% =================================================================
% GUIA RÁPIDO DE LaTeX (EXEMPLOS DIDÁTICOS)
% Como usar: inclua em article-main-v2.tex se desejar.
% =================================================================
\section{Guia Rápido de LaTeX}

\subsection{Citações e Referências}
Use \verb+\cite{chave}+ para citar (ex.: \cite{ref:exemplo}). Use \verb+\label{...}+ e \verb+\ref{...}+ para seções/figuras/tabelas.

\subsection{Figuras lado a lado (subfiguras)}
% Nota: o template SBC usa caption2; evite o pacote caption/subcaption.
% Use minipage e legendas textuais por subfigura, com uma legenda geral.
\begin{figure}[h]
  \centering
  \begin{minipage}{.48\linewidth}
    \centering
    \includegraphics[width=\linewidth]{template_diagram1.png}\\
    \textit{(a) [Subfigura A]}
    \label{fig:sub-a}
  \end{minipage}\hfill
  \begin{minipage}{.48\linewidth}
    \centering
    \includegraphics[width=\linewidth]{template_diagram2.png}\\
    \textit{(b) [Subfigura B]}
    \label{fig:sub-b}
  \end{minipage}
  \caption{[Exemplo de figuras lado a lado]}
  \label{fig:subcaption-exemplo}
\end{figure}

\subsection{Tabela Larga com tabularx}
\begin{table}[h]
  \centering
  \caption{[Tabela com colunas flexíveis]}
  \label{tab:tabularx-exemplo}
  \begin{tabularx}{\linewidth}{lXX}
    \toprule
    Campo & Descrição (coluna X) & Observações (coluna X) \\
    \midrule
    A & Texto longo que se ajusta automaticamente à largura disponível & Outra observação longa para demonstração \\
    B & Texto adicional & Mais observações \\
    \bottomrule
  \end{tabularx}
\end{table}

\subsection{Equações Numeradas e Referências}
Uma equação de exemplo em~\eqref{eq:exemplo}:
\begin{equation}
  E = m c^{2}
  \label{eq:exemplo}
\end{equation}

\subsection{Pseudocódigo / Código}
Com \texttt{listings} (não requer shell-escape):
\begin{lstlisting}[style=code, caption={Exemplo com listings}, label={lst:listings-exemplo}]
for i in 1..N {
  // faça algo
}
\end{lstlisting}

% Com minted (opcional) — requer TEXFLAGS="-shell-escape" e \usepackage{minted}
% \begin{minted}[fontsize=\footnotesize, linenos]{python}
% for i in range(10):
%     print(i)
% \end{minted}
